%% appD-terminology \dash{} Polish--English Terminology
%
% The rows of section D.3 and D.4 are \plterm{}{}{}, which is greppable on
% purpose: `check_structure.py --terms` requires every Polish rendering named
% here to occur in the prose of programs/pl. A glossary is a claim about the
% body on every line, and the failure that matters is a row naming a word the
% book does not use -- which notes/03's suggested table had two of.
%
% Section D.1's table is NOT gated and is hand-authored on purpose. Both its
% columns are literal text, because writing the contract macro there would
% print the current edition's form in BOTH columns and say nothing -- and
% because C10 forbids \tan, \cot, \arctan, \gcd and \operatorname{lcm}
% outright, rightly, since prose may never hard-code a form the other edition
% sets differently. Parity's C3 gates that both language files define the same
% macro SET; nothing checks that this table describes them correctly. So read
% it against lang/en.tex and lang/pl.tex whenever either changes.

\chapter{Polish--English Terminology}
\label{app:D}

\noindent
Two languages, one book, and one page where the choices are written down. This
appendix is for a reader of either edition: the Polish one, who needs to know
which English word a Polish term is standing in for, and the English one, who
may be reading a Polish paper or sitting in a Polish design review.

It also says where the Polish edition has \emph{not} settled on one word. That
is the harder half and it is here rather than tidied away, because a reader who
meets two words for one thing needs to be told they are one thing.

\section{The symbols the two editions set differently}
\label{sec:D-notation}
\index{notation!contract}
\index{terminology!notation contract}

The rule is the \textbf{keyboard test}: if the reader will type the token into a
computer, both editions use the code spelling; if they will only ever read it or
write it by hand, the Polish edition uses the Polish form. So the Polish edition
writes \emph{tg} for the tangent and $\tanh$ \dash{} not \emph{tgh} \dash{} for
the hyperbolic tangent, and that apparent inconsistency is the rule working
rather than an oversight.

Each row below is one macro in \code{lang/en.tex} and \code{lang/pl.tex}, so the
\emph{source} of both editions is identical and only the output differs. A macro
defined in one language file and not the other is an undefined control sequence
in exactly one of the two builds, which is why the pair is gated.

Both columns are set as literal text on purpose. Writing the contract macro here
would print the current edition's form in both columns and say nothing.

\begin{center}
\begin{tabularx}{\linewidth}{@{}llX@{}}
\toprule
\textbf{English sets} & \textbf{Polish sets} & \textbf{the macro, and what it is} \\
\midrule
\emph{tan}    & \emph{tg}     & \code{\textbackslash tg}, the tangent \\
\emph{cot}    & \emph{ctg}    & \code{\textbackslash ctg}, the cotangent \\
\emph{arctan} & \emph{arc tg} & \code{\textbackslash arctg}, the inverse tangent \\
\emph{arccot} & \emph{arc ctg} & \code{\textbackslash arcctg}, the inverse cotangent \\
\emph{gcd}    & \emph{NWD}    & \code{\textbackslash gcdop}, the greatest common divisor \\
\emph{lcm}    & \emph{NWW}    & \code{\textbackslash lcmop}, the lowest common multiple \\
\emph{span}   & \emph{lin}    & \code{\textbackslash spanof}, the span of a set of vectors \\
0.5           & 0,5           & \code{\textbackslash num} and \code{\textbackslash val}, the decimal marker \\
\bottomrule
\end{tabularx}
\end{center}

\medskip
And some that are deliberately \emph{not} on that list, because both editions
set them the same way. $\Prob$, $\Ex$, $\Var$ and $\Cov$ are written as they are
written in English, and the interval $\intcc{a}{b}$ keeps its square brackets in
both. The variance and the interval are the two a Polish reader is most likely
to be surprised by, and both are taken up in section~\ref{sec:D-unsettled}.

\section{Where the book stops to say a symbol is not what you expect}
\label{sec:D-boxes}
\index{notation!collisions}

A notation box is the book stopping mid-frame to say that a symbol or a word is
about to mean something other than what the reader brings to it. They are worth
listing in one place, because each one marks a collision somebody has already
walked into.

\begin{center}
\begin{tabularx}{\linewidth}{@{}lX@{}}
\toprule
\textbf{where} & \textbf{what it settles} \\
\midrule
\ref{prog:F02} & a letter is a name, not a mystery \\
\ref{prog:F03} & never a bare logarithm, and why the ban has exceptions \\
\ref{prog:F04} & one English word doing two jobs \\
\ref{prog:P01} & mantissa, and what the word carries \\
\ref{prog:P03} & a logarithm with no base, inside an $O$ \\
\ref{prog:P12} & the binomial coefficient, read aloud two ways \\
\ref{prog:P15} & \emph{gradient} does two jobs, a slope and a vector \\
\ref{prog:P18} & numerator layout, declared once for the book \\
\ref{prog:P23} & \emph{independent} does two jobs \\
\ref{prog:P24} & capitals and lower case; an expectation is a number; and
                 $\Var(X)$ against $D^{2}(X)$ \\
\bottomrule
\end{tabularx}
\end{center}

\medskip
The last of those is the one that matters most to a Polish reader and it is not
optional: $D^{2}(X)$ and $M(X)$ are \emph{current} Polish teaching usage rather
than a historical curiosity, so a reader who learned them and meets an
unexplained $\Var$ has reason to doubt the book rather than their schooling.
The book writes $\Var$ and $\Ex$ anyway, and the reason is the one this whole
appendix runs on: those are what the libraries and the papers a reader is going
to open write. Program~\ref{prog:P24} carries the correspondence in both
directions, so neither notation is left standing alone.

\section{Terms, English-headed}
\label{sec:D-terms}
\index{terminology!glossary}

The English term is the head, because it is what the reader will search for,
read in a paper and see in an API. Where a Polish form is in real use, the
Polish edition uses it. Where none is, the Polish edition keeps the English word
and inflects it Polish \dash{} \emph{embeddingów}, \emph{tokenizator},
\emph{benchmarku} \dash{} which is what Polish engineers say, and a calque
invented by a translator is worse than a borrowing everybody already uses.

The three tables are that rule, sorted. These the Polish edition translates.

\begin{center}
\begin{tabularx}{\linewidth}{@{}llX@{}}
\toprule
\textbf{English} & \textbf{Polish edition} & \textbf{note} \\
\midrule
\plterm{activation}{aktywacja}{Established. \emph{Funkcja aktywacji} for the
  function itself.}
\plterm{attention}{uwaga}{\emph{mechanizm uwagi}, \emph{warstwa uwagi},
  \emph{głowica uwagi}. Real Polish usage; the bare English term is not used.}
\plterm{backpropagation}{propagacja wsteczna}{Established.
  Program~\ref{prog:F12} says what it is: the chain rule applied to a
  composition of layers.}
\plterm{cross-entropy}{entropia krzyżowa}{Established.}
\plterm{evaluation set}{zbiór ewaluacyjny}{Established.}
\plterm{feature}{cecha}{Established, and used throughout.}
\plterm{gradient descent}{spadek gradientu}{Established. The book writes
  both \emph{spadek gradientu} and \emph{spadek gradientowy}, which is one
  name in two grammatical forms rather than two words.}
\plterm{layer}{warstwa}{Established.}
\plterm{loss}{strata}{Established. \emph{Funkcja straty} for the function.}
\plterm{optimiser}{optymalizator}{Established.}
\plterm{overfitting}{przeuczenie}{Established.}
\plterm{perplexity}{perpleksja}{Established. Program~\ref{prog:P29} says what
  the number counts.}
\plterm{probe}{sonda}{An ordinary Polish word with real use in this sense,
  so not a calque. Program~\ref{prog:P31}.}
\plterm{sequence}{sekwencja}{Established.}
\plterm{vocabulary}{słownik}{Established.}
\plterm{weight}{waga}{Established.}
\bottomrule
\end{tabularx}
\end{center}

\medskip
And these it borrows, because no Polish form is in real use. The
English word is kept and inflected Polish, which is what Polish engineers
say.
\begin{center}
\begin{tabularx}{\linewidth}{@{}llX@{}}
\toprule
\textbf{English} & \textbf{Polish edition} & \textbf{note} \\
\midrule
\plterm{benchmark}{benchmark}{Borrowed and inflected.}
\plterm{checkpointing}{checkpointing}{Borrowed. Trading arithmetic for bytes,
  Program~\ref{prog:P16}.}
\plterm{harness}{harness}{Borrowed and inflected; \emph{zestaw} was tried and
  collided with \emph{zbiór}.}
\plterm{hyperparameter}{hiperparametr}{Borrowed and inflected.}
\plterm{logit}{logit}{The same word in Polish statistics.}
\plterm{prompt}{prompt}{Borrowed and inflected.}
\plterm{softmax}{softmax}{Borrowed; not translated anywhere.}
\plterm{token}{token}{Borrowed and inflected. \emph{Tokenizator} for the thing
  that produces them.}
\plterm{transformer}{transformer}{Borrowed. Never \emph{transformator}, which
  is an electrical device.}
\bottomrule
\end{tabularx}
\end{center}

\medskip
These need more than a gloss. In each, one word is doing two jobs, or
the two editions have settled on different words, and a reader who takes the
first reading will be wrong somewhere else in the book.
\begin{center}
\begin{tabularx}{\linewidth}{@{}llX@{}}
\toprule
\textbf{English} & \textbf{Polish edition} & \textbf{note} \\
\midrule
\plterm{head}{głowica}{Of attention. \emph{Głowa} is the anatomical one.}
\plterm{inference}{wnioskowanie}{Two jobs, and the Polish carries the English
  collision rather than resolving it: statistical inference,
  Programs~\ref{prog:P27} and \ref{prog:P28}, and running a trained model, as
  in \emph{przy wnioskowaniu} at Program~\ref{prog:F02}.}
\plterm{normalisation}{normalizacja}{Two jobs, in Polish as in English:
  dividing a vector by its length, Program~\ref{prog:F09}, and layer
  normalisation, Program~\ref{prog:P32}.}
\plterm{precision}{precyzja}{Two jobs, and the collision is in both languages:
  the classification metric of Program~\ref{prog:P23}, and the width of a
  floating-point significand throughout Program~\ref{prog:P01}.}
\plterm{recall}{czułość}{Established. It sits beside \emph{precyzja} on the
  same classification report and conditions on the other population.}
\plterm{residual connection}{rezydualny}{\emph{Połączenie rezydualne},
  \emph{strumień rezydualny}. Program~\ref{prog:P32}.}
\plterm{residual of a fit}{reszta}{A different object, and Polish gives it a
  different word: what a least-squares fit leaves over,
  Program~\ref{prog:P08}. English uses one word for both.}
\plterm{significand}{mantysa}{The two editions disagree here and both are
  right. The standard says \emph{significand} for the stored bits and reserves
  \emph{mantissa} for the fractional part of a logarithm, so the English
  edition writes \emph{significand}; Polish practice says \emph{mantysa} for
  both objects, so the Polish edition does too. Program~\ref{prog:P01} carries
  a notation box in each edition, reaching the opposite conclusion in each.}
\bottomrule
\end{tabularx}
\end{center}

\section{Where Polish usage has not settled}
\label{sec:D-unsettled}
\index{terminology!unsettled}

Four terms in this book are rendered more than one way in the Polish edition.
They are listed rather than quietly picked, because a reader who meets the
second word needs to know it is the first one.

\begin{center}
\begin{tabularx}{\linewidth}{@{}llX@{}}
\toprule
\textbf{English} & \textbf{Polish edition} & \textbf{note} \\
\midrule
\plterm{batch}{partia, wsad, batch}{One object, and each of the renderings
  beside it has a job. \emph{Partia} is the base: it is the commonest, and the
  only one that compounds, which is where \emph{mikropartia} comes from.
  \emph{Wsad} survives for the adjective \emph{wsadowy}, which \emph{partia}
  cannot form. \emph{Batch} is the borrowing and stays where the reader would
  type it into a search box. How many there are is deliberately not counted
  here: it was once, and a fourth had appeared underneath it.}
\plterm{embedding}{zanurzenie, embedding}{This book's own translation rule says
  to keep the English word where the Polish form is a calque nobody searches
  for, and the Polish edition mostly uses \emph{zanurzenie} anyway. The reason
  is the verb: \emph{zanurzać} has no borrowed form \dash{} nobody says
  \emph{embedduje} \dash{} so the noun followed it. \emph{Embedding} is the
  term a Polish engineer will type into a search box.}
\plterm{learning rate}{współczynnik uczenia, krok uczenia}{Both are in real
  use. \emph{Krok} reads as \emph{step}, which is what it multiplies.}
\plterm{sample}{próba, próbka}{\emph{Próba} is the statistical term and is what
  the book writes in the collocations \dash{} \emph{wariancja z próby},
  \emph{średnia z próby}, Program~\ref{prog:P26}. \emph{Próbka} is what it
  writes where the sample is a concrete thing in hand, as in
  Program~\ref{prog:P27}'s bootstrap. \emph{Próbkowanie} is the operation
  either way.}
\bottomrule
\end{tabularx}
\end{center}

\medskip
\begin{note}
The sharpest cost of an unsettled word is a cross-reference that renames the
thing it points at. Program~\ref{prog:P32}'s table of what this book has not
measured sends the reader to Program~\ref{prog:P08} and Program~\ref{prog:P11}
for the spectrum of a real embedding matrix; the Polish edition writes
\emph{macierz zanurzeń} there, and so does Program~\ref{prog:P11}, while
Program~\ref{prog:P08} writes \emph{macierz embeddingów}. That one cannot be
repaired at the pointer, because the row names two programs that disagree with
each other, and it is the argument for settling \emph{embedding} before the
other two.
\end{note}

\medskip
Two more entries are worth naming because they are the opposite case \dash{}
places where the Polish edition is settled and the reader may expect otherwise.
$\intcc{a}{b}$ keeps square brackets rather than the angle brackets Polish
schooling uses, because ISO 80000-2 and every paper the reader is being trained
to open write them that way. No program stops to explain that, so
Appendix~\ref{app:B} is where it is recorded rather than the first interval a
reader meets. And the hyperbolic tangent stays $\tanh$ rather than \emph{tgh},
which is the keyboard test of section~\ref{sec:D-notation} deciding against the
Polish form because the reader will type this one.

\section{Two rules nothing can check}
\label{sec:D-untestable}
\index{terminology!translator rules}

The two editions are compared against each other line by line, and two of the
translator's rules sit outside everything that comparison can see.

\textbf{A listing's comments stay English.} The reader is being trained to read
English code, and a comment translated into Polish is the one line of a session
they could not paste into a terminal and search for. Most of the book's
listings are console transcripts generated once by \code{code/} and pulled into
both editions from the same file, so their comments are English by construction
rather than by anybody's discipline; a listing authored inline is held to the
rule by hand.

\textbf{A constant's spelling is never translated.} $e$ is $e$ and $\pi$ is
$\pi$ in both editions. They are names rather than words, so translating one
would be translating an identifier \dash{} which is the keyboard test of
section~\ref{sec:D-notation} again, applied to a symbol the reader will type
rather than to a function they will call.
